\chapter{Kunihiko Kodaira --- Fields Medal 1954}\label{ch1954kunihikokodaira:kunihiko-kodaira-fields-medal-1954-a-comprehensive-report}

\begin{center}\rule{0.5\linewidth}{0.5pt}\end{center}

\textbf{Prepared:} July 2026 \textbf{Subject:} Kunihiko Kodaira (b.
March 16, 1915, Tokyo, Japan) \textbf{Position:} Professor, Stanford
University; Princeton University \textbf{Doctoral Advisor:} Shokichi
Iyanaga (University of Tokyo) \textbf{Fields Medal Citation:}
\emph{``Achieved major results in the theory of harmonic integrals and
numerous applications to Kählerian and more specifically to algebraic
varieties. He demonstrated, by sheaf cohomology, that such varieties are
Hodge manifolds.''}

\begin{center}\rule{0.5\linewidth}{0.5pt}\end{center}

\section{1. Introduction}\label{ch1954kunihikokodaira:introduction}

Kunihiko Kodaira was awarded the Fields Medal in 1954 for his significant
contributions to complex algebraic geometry, particularly the theory of
harmonic integrals, the Kodaira vanishing theorem, the classification of
complex surfaces, and the deformation theory of complex structures. He
was the first Japanese mathematician to receive the Fields Medal, and
his work laid the foundation for much of modern algebraic geometry.

Kodaira's approach was unique in its synthesis of analytic methods
(Hodge theory, elliptic PDEs) and algebraic methods (sheaf cohomology,
the Riemann--Roch theorem). He was able to prove results of notable
depth using the interaction between these two perspectives.

His principal legacy can be seen in several directions: the Kodaira
vanishing theorem is one of the most used tools in algebraic geometry;
the Kodaira embedding theorem characterizes projective manifolds
geometrically; the classification of surfaces remains the definitive
organizational framework for complex surfaces; and the Kodaira--Spencer
deformation theory provides the foundation for the study of moduli
spaces.

Kodaira's major contributions include:

\begin{itemize}
\tightlist
\item
 The \textbf{Kodaira vanishing theorem}, a fundamental result in
 algebraic geometry that guarantees the vanishing of certain sheaf
 cohomology groups on projective varieties, essential for the Kodaira
 embedding theorem and the minimal model program.
\item
 The \textbf{Kodaira embedding theorem}, characterizing projective
 manifolds as those carrying a Hodge metric (a Kähler metric with
 integral cohomology class).
\item
 The \textbf{Kodaira--Spencer deformation theory}, establishing the
 foundational framework for studying variations of complex structures.
\item
 The \textbf{classification of compact complex surfaces}, organizing
 all such surfaces into seven classes according to their Kodaira
 dimension.
\item
 The \textbf{Kodaira dimension} \(\kappa(X)\), a birational invariant
 that measures the growth of the canonical ring.
\item
 The \textbf{Kodaira--Serre duality}, relating sheaf cohomology groups
 on complex manifolds (a form of Poincaré duality for complex
 manifolds).
\item
 The \textbf{Riemann--Roch theorem for complex manifolds} (the
 Hirzebruch--Riemann--Roch theorem), proved using harmonic integrals.
\item
 The \textbf{theory of elliptic surfaces} and their classification.
\end{itemize}

This report provides a comprehensive overview of Kodaira's life, work,
and lasting impact on algebraic geometry and complex analysis.

\begin{center}\rule{0.5\linewidth}{0.5pt}\end{center}

\section{2. Biographical Background}\label{ch1954kunihikokodaira:biographical-background}

\subsection{2.1 Early Life and
Education}\label{ch1954kunihikokodaira:early-life-and-education}

Kunihiko Kodaira was born in Tokyo on March 16, 1915. He entered the
University of Tokyo in 1935, studying under the eminent mathematician
Teiji Takagi (known for Takagi's theorem on class field theory) and
Shokichi Iyanaga. In his first year, he took Takagi's ``Introduction to
Analysis'' course, which sparked his interest in analysis. He graduated
with a Bachelor of Science in mathematics in 1938 and then earned a
second bachelor's degree in physics in 1941, already having ten
published papers by then. During his physics studies, he became close to
Iyanaga's family and eventually married Iyanaga's sister Seiko, a
violinist, in 1943.

\subsection{2.2 Academic Career}\label{ch1954kunihikokodaira:academic-career}

\begin{itemize}
\tightlist
\item
 \textbf{1941--1944:} Lecturer, University of Tokyo (Physics
 Department).
\item
 \textbf{1944--1949:} Assistant Professor, University of Tokyo.
\item
 \textbf{1949:} PhD, University of Tokyo (dissertation: ``Harmonic
 Fields in Riemannian Manifolds'').
\item
 \textbf{1949--1952:} Member, Institute for Advanced Study, Princeton.
\item
 \textbf{1952--1955:} Associate Professor, Princeton University.
\item
 \textbf{1954:} Fields Medal awarded at the ICM in Amsterdam.
\item
 \textbf{1955--1962:} Professor, Princeton University.
\item
 \textbf{1962--1965:} Professor, Johns Hopkins University.
\item
 \textbf{1965--1975:} Professor, Stanford University.
\item
 \textbf{1975--1984:} Professor, Gakushuin University, Tokyo.
\item
 \textbf{1997:} Died at age 82 in Kōfu, Japan.
\end{itemize}

\subsection{2.3 Doctoral Work}\label{ch1954kunihikokodaira:doctoral-work}

Kodaira's PhD dissertation on harmonic fields in Riemannian manifolds,
completed under Shokichi Iyanaga, already contained the seeds of his
life's work. He developed Hodge theory in a new direction, generalizing
Hodge's results to manifolds with singularities. Weyl recognized the
importance of this work and invited Kodaira to the Institute for
Advanced Study in Princeton in 1949.

\subsection{2.4 Personal Traits and Musical
Interests}\label{ch1954kunihikokodaira:personal-traits-and-musical-interests}

Kodaira was an accomplished pianist who often played duets with his wife
Seiko (a violinist). His colleagues noted that his mathematical
creativity was intertwined with his musical sensibility. He once
remarked that the structure of a Beethoven symphony and the structure of
a mathematical proof were not so different.

\begin{center}\rule{0.5\linewidth}{0.5pt}\end{center}

\section{3. Harmonic Integrals and Hodge
Theory}\label{ch1954kunihikokodaira:harmonic-integrals-and-hodge-theory}

\subsection{3.1 Hodge Theory}\label{ch1954kunihikokodaira:hodge-theory}

\textbf{Hodge theory} (developed by W. V. D. Hodge in the 1930s) studies
harmonic differential forms on compact Riemannian manifolds. On a
compact Kähler manifold, the space of harmonic \(p\)-forms decomposes
into types \((r, s)\) with \(r + s = p\).

\subsection{3.2 Kodaira's Contributions}\label{ch1954kunihikokodaira:kodairas-contributions}

Kodaira made several fundamental contributions to Hodge theory: - He
generalized Hodge's theory of harmonic integrals to manifolds with
boundary. - He developed the theory of \textbf{harmonic forms with
singularities}. - He used Hodge theory to prove the \textbf{Kodaira
embedding theorem}, showing that a Hodge manifold is algebraic.

\subsection{3.3 Harmonic Forms with
Singularities}\label{ch1954kunihikokodaira:harmonic-forms-with-singularities}

Kodaira's work on harmonic forms with singularities allowed him to study
the behavior of harmonic forms near a divisor on a complex manifold. His
approach was to consider harmonic forms on a manifold \(X \setminus D\)
where \(D\) is a divisor, and study their asymptotic behavior near
\(D\). This was essential for: - Proving the \textbf{vanishing theorem}.
- Constructing \textbf{meromorphic functions} on algebraic varieties. -
Generalizing the \textbf{Riemann--Roch theorem} to higher dimensions.

\subsection{3.4 The Weitzenböck
Formula}\label{ch1954kunihikokodaira:the-weitzenbuxf6ck-formula}

The \textbf{Weitzenböck formula} (also called the Bochner--Kodaira
formula) expresses the Hodge Laplacian in terms of the rough Laplacian
and a curvature term:

\[\Delta = \nabla^* \nabla + \mathcal{R}\]

where \(\mathcal{R}\) is a curvature operator. Kodaira used this formula
to prove vanishing theorems by showing that when the curvature term is
positive, the Laplacian has a positive lower bound, forcing cohomology
groups to vanish.

\begin{center}\rule{0.5\linewidth}{0.5pt}\end{center}

\section{4. The Kodaira Vanishing
Theorem}\label{ch1954kunihikokodaira:the-kodaira-vanishing-theorem}

\subsection{4.1 Statement}\label{ch1954kunihikokodaira:statement}

The \textbf{Kodaira vanishing theorem} is one of the important
results in algebraic geometry. It states:

\begin{quote}
\textbf{Kodaira Vanishing Theorem:} Let \(X\) be a compact Kähler
manifold of dimension \(n\), and let \(L\) be a positive line bundle on
\(X\) (i.e., a line bundle whose first Chern class is represented by a
Kähler form). Then:
\[H^q(X, K_X \otimes L) = 0 \quad \text{for all } q > 0\] where \(K_X\)
is the canonical bundle of \(X\).
\end{quote}

\subsection{4.2 Equivalent Formulation}\label{ch1954kunihikokodaira:equivalent-formulation}

Equivalently, for a positive line bundle \(L\):

\[H^q(X, L^{-1}) = 0 \quad \text{for all } q < n\]

This formulation shows that the cohomology of line bundles with
sufficiently negative curvature vanishes in all but the top degree.

\subsection{4.3 Proof via Hodge Theory}\label{ch1954kunihikokodaira:proof-via-hodge-theory}

The proof uses Hodge theory and the \textbf{Bochner--Kodaira--Nakano
identity}:

\[\bar{\partial} \bar{\partial}^* + \bar{\partial}^* \bar{\partial} + \bar{\partial}_\theta \bar{\partial}_\theta^* + \bar{\partial}_\theta^* \bar{\partial}_\theta = \sqrt{-1} [\Theta(L), \Lambda]\]

where \(\Theta(L)\) is the curvature form of \(L\) and \(\Lambda\) is
the adjoint of multiplication by the Kähler form. From this identity,
one obtains a positive lower bound for the Laplacian on \((n, q)\)-forms
with values in \(L\).

\subsection{4.4 Generalizations}\label{ch1954kunihikokodaira:generalizations}

The Kodaira vanishing theorem has been generalized in many directions: -
\textbf{Kawamata--Viehweg vanishing} for nef and big line bundles on
singular varieties. - \textbf{Kollár vanishing} for higher direct images
of dualizing sheaves. - \textbf{Nadel vanishing} for multiplier ideal
sheaves. - \textbf{Mckernan vanishing} for log canonical pairs.

\begin{center}\rule{0.5\linewidth}{0.5pt}\end{center}

\section{5. The Kodaira Embedding
Theorem}\label{ch1954kunihikokodaira:the-kodaira-embedding-theorem}

\subsection{5.1 Statement}\label{ch1954kunihikokodaira:statement-1}

The \textbf{Kodaira embedding theorem} characterizes projective
algebraic manifolds among complex manifolds:

\begin{quote}
\textbf{Kodaira Embedding Theorem:} A compact complex manifold \(X\) is
projective algebraic (i.e., can be embedded in projective space) if and
only if it carries a \textbf{Hodge metric} --- a Kähler metric whose
cohomology class lies in \(H^2(X, \mathbb{Z})\).
\end{quote}

\subsection{5.2 Hodge Manifolds}\label{ch1954kunihikokodaira:hodge-manifolds}

A \textbf{Hodge manifold} is a compact Kähler manifold whose Kähler form
\(\omega\) is integral: \([\omega] \in H^2(X, \mathbb{Z})\). The Kodaira
embedding theorem states that every Hodge manifold is projective.

\subsection{5.3 Proof Outline}\label{ch1954kunihikokodaira:proof-outline}

The proof proceeds in several steps: 1. The integrality condition
implies that \(\omega\) is the curvature of a line bundle \(L\) (the
\textbf{Lefschetz theorem on \((1,1)\)-classes}). 2. By the Kodaira
vanishing theorem, \(H^q(X, L^k \otimes K_X^{-1}) = 0\) for \(q > 0\)
and \(k\) sufficiently large. 3. The Riemann--Roch theorem computes
\(\dim H^0(X, L^k)\) grows as \(k^n\), giving enough sections. 4. The
sections of \(L^k\) define a holomorphic embedding of \(X\) into
projective space.

\subsection{5.4 The Chow--Kodaira
Theorem}\label{ch1954kunihikokodaira:the-chowkodaira-theorem}

A related result is the \textbf{Chow--Kodaira theorem} (proved by
Kodaira and Wei-Liang Chow): a compact complex manifold that is a
submanifold of projective space is algebraic in the sense that it is the
zero set of homogeneous polynomials. This result shows that the analytic
and algebraic definitions of projective variety coincide for smooth
varieties.

\subsection{5.5 Consequences}\label{ch1954kunihikokodaira:consequences}

The embedding theorem has numerous consequences: - Every compact Riemann
surface is projective (the classical fact). - Every complex torus with a
Riemann form is an abelian variety. - The classification of complex
surfaces is intimately related to projectivity. - The \textbf{Moĭshezon
theorem}: a compact complex manifold of dimension \(n\) with \(n\)
algebraically independent meromorphic functions is bimeromorphic to a
projective variety.

\begin{center}\rule{0.5\linewidth}{0.5pt}\end{center}

\section{6. The Kodaira Dimension}\label{ch1954kunihikokodaira:the-kodaira-dimension}

\subsection{6.1 Definition}\label{ch1954kunihikokodaira:definition}

The \textbf{Kodaira dimension} \(\kappa(X)\) of a compact complex
manifold \(X\) is defined in terms of the growth of the
\textbf{canonical ring}:

\[R(X) = \bigoplus_{m \geq 0} H^0(X, K_X^m)\]

Let \(P_m = \dim H^0(X, K_X^m)\) be the \(m\)-th \textbf{plurigenus}.
The Kodaira dimension is:

\[\kappa(X) =
\begin{cases}
-\infty & \text{if } P_m = 0 \text{ for all } m \geq 1 \\
\limsup_{m \to \infty} \frac{\log P_m}{\log m} & \text{otherwise}
\end{cases}\]

\subsection{6.2 Values}\label{ch1954kunihikokodaira:values}

The Kodaira dimension takes values
\(\kappa \in \{-\infty, 0, 1, 2, \ldots, n\}\) where \(n = \dim X\): -
\(\kappa(X) = -\infty\): \(X\) is of \textbf{general type} if
\(\kappa = n\), or \textbf{uniruled} if \(\kappa = -\infty\). -
\(\kappa(X) = 0\): The canonical bundle is torsion (Calabi--Yau
manifolds, abelian varieties). - \(\kappa(X) = n\): \(X\) is of
\textbf{general type}.

\subsection{6.3 The Plurigenera}\label{ch1954kunihikokodaira:the-plurigenera}

The \(m\)-th \textbf{plurigenus} \(P_m = \dim H^0(X, K_X^m)\) is a
birational invariant. For surfaces, the sequence \((P_1, P_2, \ldots)\)
determines the Kodaira dimension: - \(\kappa = -\infty\): \(P_m = 0\)
for all \(m\). - \(\kappa = 0\): \(P_m = 0\) or \(1\) for all \(m\), and
\(P_{12} = 1\). - \(\kappa = 1\): \(P_m \sim c m\) for some \(c > 0\). -
\(\kappa = 2\): \(P_m \sim c m^2\) for some \(c > 0\).

\subsection{6.4 Iitaka's Conjecture}\label{ch1954kunihikokodaira:iitakas-conjecture}

The \textbf{Iitaka conjecture} (also called the \(C_{n,m}\) conjecture)
relates the Kodaira dimension of a fibration to that of its base and
fiber. It states that for a smooth fibration \(f: X \to Y\):

\[\kappa(X) \geq \kappa(F) + \kappa(Y)\]

where \(F\) is the generic fiber. This has been proved in many cases by
Viehweg, Kawamata, and others.

\begin{center}\rule{0.5\linewidth}{0.5pt}\end{center}

\section{7. The Classification of Compact Complex
Surfaces}\label{ch1954kunihikokodaira:the-classification-of-compact-complex-surfaces}

\subsection{7.1 Kodaira's Classification}\label{ch1954kunihikokodaira:kodairas-classification}

Kodaira classified all compact complex surfaces (complex manifolds of
dimension 2) into seven classes:

\begin{enumerate}
\def\labelenumi{\arabic{enumi}.}
\tightlist
\item
 \textbf{Rational surfaces} (\(\kappa = -\infty\)): \(\mathbb{P}^2\),
 ruled surfaces over \(\mathbb{P}^1\).
\item
 \textbf{Ruled surfaces} (\(\kappa = -\infty\)):
 \(\mathbb{P}^1\)-bundles over curves of genus \(g \geq 1\).
\item
 \textbf{Surfaces of class VII} (\(\kappa = -\infty\)): Non-Kähler
 surfaces with \(b_1 = 1\).
\item
 \textbf{K3 surfaces} (\(\kappa = 0\)): Simply connected surfaces with
 trivial canonical bundle.
\item
 \textbf{Enriques surfaces} (\(\kappa = 0\)): Quotients of K3 surfaces
 by a fixed-point-free involution.
\item
 \textbf{Tori} (\(\kappa = 0\)): Complex tori of dimension 2.
\item
 \textbf{Hyperelliptic surfaces} (\(\kappa = 0\)): Quotients of tori by
 finite groups.
\item
 \textbf{Properly elliptic surfaces} (\(\kappa = 1\)): Surfaces with an
 elliptic fibration.
\item
 \textbf{Surfaces of general type} (\(\kappa = 2\)): All other
 surfaces.
\end{enumerate}

\subsection{7.2 The Classification
Method}\label{ch1954kunihikokodaira:the-classification-method}

The classification uses: - The \textbf{Kodaira dimension} as the primary
invariant. - The \textbf{Enriques--Kodaira classification} of surfaces
by their invariants \(p_g\) (geometric genus), \(q\) (irregularity), and
\(b_1\) (first Betti number). - The \textbf{Kähler condition}: surfaces
with even first Betti number are Kähler (a theorem of Kodaira). - The
\textbf{minimal model program}: every surface has a minimal model with
no \((-1)\)-curves.

\subsection{7.3 The Minimal Model}\label{ch1954kunihikokodaira:the-minimal-model}

A \textbf{minimal surface} is one with no \((-1)\)-curves (rational
curves with self-intersection \(-1\)). Every surface can be obtained
from a minimal surface by a sequence of blow-ups. Kodaira showed that
minimal surfaces are classified by their Kodaira dimension.

\subsection{7.4 The Classification
Table}\label{ch1954kunihikokodaira:the-classification-table}

The complete Enriques--Kodaira classification can be summarized as:

{\def\LTcaptype{none} % do not increment counter
\begin{longtable}[]{@{}lllll@{}}
\toprule\noalign{}
\(\kappa\) & Type & \(b_1\) & \(p_g\) & \(q\) \\
\midrule\noalign{}
\endhead
\bottomrule\noalign{}
\endlastfoot
\(-\infty\) & Rational & 0 & 0 & 0 \\
\(-\infty\) & Ruled (\(g \geq 1\)) & \(2g\) & 0 & \(g\) \\
\(-\infty\) & Class VII & 1 & 0 & 0 \\
0 & K3 & 0 & 1 & 0 \\
0 & Enriques & 0 & 0 & 0 \\
0 & Torus & 4 & 1 & 2 \\
0 & Hyperelliptic & 2 & 0 & 1 \\
1 & Properly elliptic & \(\leq 2\) & varies & varies \\
2 & General type & varies & varies & varies \\
\end{longtable}
}

\subsection{7.5 The Betti Number
Criterion}\label{ch1954kunihikokodaira:the-betti-number-criterion}

Kodaira proved a fundamental result: a compact complex surface is Kähler
if and only if \(b_1\) is even. This criterion is essential for the
classification, as surfaces with \(b_1\) odd (class VII surfaces) are
non-Kähler and have special properties.

\subsection{7.6 The Classification of Class VII
Surfaces}\label{ch1954kunihikokodaira:the-classification-of-class-vii-surfaces}

Class VII surfaces (\(\kappa = -\infty\), \(b_1 = 1\)) were completely
classified by Bogomolov, Dloussky, and others. They include: -
\textbf{Hopf surfaces:} Quotients of \(\mathbb{C}^2 \setminus \{0\}\) by
a contraction \((z,w) \mapsto (\alpha z, \beta w)\) with
\(|\alpha|, |\beta| < 1\). - \textbf{Inoue surfaces:} Certain surfaces
with no meromorphic functions. - \textbf{Kato surfaces:} Surfaces with a
cycle of rational curves.

\begin{center}\rule{0.5\linewidth}{0.5pt}\end{center}

\section{8. Deformation Theory of Complex
Structures}\label{ch1954kunihikokodaira:deformation-theory-of-complex-structures}

\subsection{8.1 The Kodaira--Spencer Map}\label{ch1954kunihikokodaira:the-kodairaspencer-map}

The \textbf{Kodaira--Spencer deformation theory}, developed jointly with
Donald C. Spencer, studies how the complex structure of a manifold
varies in families. The key object is the \textbf{Kodaira--Spencer map}:

\[\rho: T_0 B \to H^1(X_0, \Theta_{X_0})\]

where \(B\) is the parameter space of the family, \(T_0 B\) is the
tangent space at the base point, and \(\Theta_{X_0}\) is the sheaf of
holomorphic vector fields on \(X_0\).

\subsection{8.2 The Kodaira--Spencer Map in
Families}\label{ch1954kunihikokodaira:the-kodairaspencer-map-in-families}

For a family \(\pi: \mathcal{X} \to B\) with fiber
\(X_t = \pi^{-1}(t)\), the Kodaira--Spencer map assigns to a tangent
vector \(v \in T_b B\) the corresponding first-order deformation. The
image \(\rho(v)\) is a cohomology class in \(H^1(X_b, \Theta_{X_b})\).

\subsection{8.3 Obstructed and Unobstructed
Deformations}\label{ch1954kunihikokodaira:obstructed-and-unobstructed-deformations}

A deformation is \textbf{unobstructed} if there are no obstructions to
extending it to higher order; this happens when the
\textbf{Kodaira--Spencer obstruction class} in
\(H^2(X_b, \Theta_{X_b})\) vanishes.

\subsection{8.4 Kuranishi Theory}\label{ch1954kunihikokodaira:kuranishi-theory}

The \textbf{Kuranishi theorem} states that the space of deformations of
a compact complex manifold is locally isomorphic to the zero set of a
holomorphic map (the \textbf{Kuranishi map}):

\[\kappa: H^1(X, \Theta_X) \to H^2(X, \Theta_X)\]

When \(H^2(X, \Theta_X) = 0\), the deformations are unobstructed and the
Kuranishi space is smooth.

\subsection{8.5 The Kodaira--Spencer Map in
Coordinates}\label{ch1954kunihikokodaira:the-kodairaspencer-map-in-coordinates}

In local coordinates, the Kodaira--Spencer map can be described
explicitly. Given a family \(\pi: \mathcal{X} \to B\) with coordinate
\(t\) on \(B\), the complex structure on \(X_t\) is described by a
Beltrami differential \(\mu(t) \in H^1(X_0, \Theta_{X_0})\). The
Kodaira--Spencer map is:

\[\rho\left(\frac{\partial}{\partial t}\right) = \left.\frac{d\mu}{dt}\right|_{t=0}\]

This connects deformation theory to the theory of quasiconformal
mappings and Beltrami equations.

\subsection{8.6 The Maurer--Cartan
Equation}\label{ch1954kunihikokodaira:the-maurercartan-equation}

The integrability condition for a deformation is governed by the
\textbf{Maurer--Cartan equation}. A 1-form
\(\varphi \in \Omega^{0,1}(X, \Theta_X)\) defines a deformation if and
only if:

\[\bar{\partial}\varphi + \frac{1}{2}[\varphi, \varphi] = 0\]

This equation, derived by Kodaira and Spencer, shows that the
deformation theory of complex structures is controlled by a differential
graded Lie algebra, a perspective that led to the development of
\textbf{derived deformation theory}.

\subsection{8.7 Applications}\label{ch1954kunihikokodaira:applications}

Deformation theory has been applied to: - The study of \textbf{moduli
spaces} of complex structures. - The classification of \textbf{Fano
varieties} and \textbf{Calabi--Yau manifolds}. - The construction of
\textbf{mirror manifolds} in string theory. - The \textbf{Torelli
problem} for K3 surfaces.

\begin{center}\rule{0.5\linewidth}{0.5pt}\end{center}

\section{9. The Riemann--Roch Theorem}\label{ch1954kunihikokodaira:the-riemannroch-theorem}

\subsection{9.1 Classical Riemann--Roch}\label{ch1954kunihikokodaira:classical-riemannroch}

For a compact Riemann surface \(C\) of genus \(g\), the
\textbf{Riemann--Roch theorem} states:

\[\dim H^0(C, L) - \dim H^1(C, L) = \deg(L) + 1 - g\]

where \(L\) is a line bundle.

\subsection{9.2 The Hirzebruch--Riemann--Roch
Theorem}\label{ch1954kunihikokodaira:the-hirzebruchriemannroch-theorem}

Kodaira contributed to the generalization of the Riemann--Roch theorem
to higher dimensions. The \textbf{Hirzebruch--Riemann--Roch theorem}
(HRR) states that for a vector bundle \(E\) on a compact complex
manifold \(X\):

\[\chi(X, E) = \int_X \text{ch}(E) \cdot \text{td}(X)\]

where \(\chi(X, E) = \sum_{i=0}^{\dim X} (-1)^i \dim H^i(X, E)\),
\(\text{ch}(E)\) is the Chern character, and \(\text{td}(X)\) is the
Todd class.

\subsection{9.3 Kodaira's Contribution}\label{ch1954kunihikokodaira:kodairas-contribution}

Kodaira proved the HRR for algebraic surfaces using Hodge theory and the
theory of harmonic integrals. His work provided the crucial analytic
foundation for the general theorem, later proved by Hirzebruch using the
cobordism theory of Thom.

\subsection{9.4 The Hirzebruch Signature
Theorem}\label{ch1954kunihikokodaira:the-hirzebruch-signature-theorem}

A related result is the \textbf{Hirzebruch signature theorem}:

\[\tau(X) = \int_X L(X)\]

where \(\tau(X)\) is the signature of the intersection form on
\(H^{2k}(X, \mathbb{R})\) and \(L(X)\) is the Hirzebruch L-genus.

\begin{center}\rule{0.5\linewidth}{0.5pt}\end{center}

\section{10. The Kodaira--Spencer--Nirenberg
Theory}\label{ch1954kunihikokodaira:the-kodairaspencernirenberg-theory}

\subsection{10.1 The Kähler Condition and
Deformations}\label{ch1954kunihikokodaira:the-kuxe4hler-condition-and-deformations}

Kodaira, Spencer, and Louis Nirenberg studied the stability of the
Kähler condition under deformations. They proved:

\begin{quote}
\textbf{Kodaira--Spencer--Nirenberg Theorem:} A small deformation of a
Kähler manifold remains Kähler.
\end{quote}

This result shows that the Kähler condition is \textbf{open} in the
moduli space of complex structures.

\subsection{10.2 The Stability Theorem}\label{ch1954kunihikokodaira:the-stability-theorem}

The proof uses the theory of elliptic PDEs and the
\textbf{Kodaira--Spencer--Nirenberg stability theorem}, which shows that
if \(X_0\) is Kähler, then for any small deformation \(X_t\), one can
construct a Kähler metric on \(X_t\) by perturbing the Kähler form.

\subsection{10.3 The Kähler Cone}\label{ch1954kunihikokodaira:the-kuxe4hler-cone}

The space of Kähler classes on a compact complex manifold \(X\) is an
open cone in \(H^{1,1}(X, \mathbb{R})\) called the \textbf{Kähler cone}
\(\mathcal{K}_X\). Kodaira's embedding theorem characterizes the
intersection \(\mathcal{K}_X \cap H^2(X, \mathbb{Z})\) as the set of
classes of Hodge metrics.

\subsection{10.4 The Nakai--Moishezon
Criterion}\label{ch1954kunihikokodaira:the-nakaimoishezon-criterion}

The \textbf{Nakai--Moishezon criterion} gives a numerical
characterization of ampleness: a line bundle \(L\) on a projective
variety \(X\) is ample if and only if \(L^{\dim Y} \cdot Y > 0\) for
every subvariety \(Y \subset X\). This criterion, combined with
Kodaira's results, provides a complete characterization of projective
varieties.

\begin{center}\rule{0.5\linewidth}{0.5pt}\end{center}

\section{11. The Kodaira--Nirenberg--Spencer
Theorem}\label{ch1954kunihikokodaira:the-kodairanirenbergspencer-theorem}

\subsection{11.1 Statement}\label{ch1954kunihikokodaira:statement-2}

The \textbf{Kodaira--Nirenberg--Spencer theorem} (also called the
\textbf{Kodaira stability theorem}) states that small deformations of a
projective algebraic manifold remain projective.

\begin{quote}
\textbf{KNS Theorem:} If \(X\) is a projective algebraic manifold, then
any sufficiently small deformation of \(X\) is also projective.
\end{quote}

\subsection{11.2 Consequences}\label{ch1954kunihikokodaira:consequences-1}

This theorem shows that projectivity is an open property in families of
complex manifolds. Together with the Kodaira embedding theorem, it
provides a complete picture of the relationship between Kähler and
projective manifolds.

\begin{center}\rule{0.5\linewidth}{0.5pt}\end{center}

\section{12. Elliptic Surfaces}\label{ch1954kunihikokodaira:elliptic-surfaces}

\subsection{12.1 Definition}\label{ch1954kunihikokodaira:definition-1}

An \textbf{elliptic surface} is a compact complex surface \(X\) equipped
with a holomorphic map \(\pi: X \to C\) to a curve \(C\) such that the
generic fiber is an elliptic curve (a torus of dimension 1).

\subsection{12.2 Kodaira's Classification of Elliptic
Surfaces}\label{ch1954kunihikokodaira:kodairas-classification-of-elliptic-surfaces}

Kodaira classified elliptic surfaces according to the types of singular
fibers that can occur. The possible singular fibers are: - \(I_n\)
(\(n \geq 1\)): A cycle of \(n\) rational curves (a ``necklace''). -
\(I_n^*\) (\(n \geq 0\)): A more complicated configuration with a
\((-2)\)-curve. - \(II, III, IV\): Fibers of types corresponding to the
Kodaira--Néron classification. - \(II^*, III^*, IV^*\): Dual fibers.

\subsection{12.3 The Canonical Bundle
Formula}\label{ch1954kunihikokodaira:the-canonical-bundle-formula}

For an elliptic surface \(\pi: X \to C\), Kodaira gave a formula for the
canonical bundle:

\[K_X = \pi^*(K_C \otimes L) \otimes \mathcal{O}_X(\sum a_i F_i)\]

where \(F_i\) are the multiple fibers and \(L\) is a line bundle on
\(C\) whose degree is related to the Euler characteristic of \(X\).

\subsection{12.4 The Canonical Class
Formula}\label{ch1954kunihikokodaira:the-canonical-class-formula}

For an elliptic surface \(\pi: X \to C\), the canonical class satisfies:

\[K_X = \pi^*(K_C + L) + \sum_i (m_i - 1)F_i\]

where \(F_i\) are the multiple fibers (of multiplicity \(m_i\)) and
\(L\) is a line bundle on \(C\) of degree \(\chi(\mathcal{O}_X)\). This
formula is essential for: - Computing the Kodaira dimension of elliptic
surfaces. - Determining when an elliptic surface is of general type.

\subsection{12.5 Applications}\label{ch1954kunihikokodaira:applications-1}

Elliptic surfaces are important in: - \textbf{String theory:} They are
used in F-theory compactifications. - \textbf{Arithmetic geometry:}
Elliptic curves over function fields. - \textbf{Mirror symmetry:}
Elliptic surfaces provide examples of mirror pairs. -
\textbf{Diophantine geometry:} The study of rational points on elliptic
curves over function fields.

\begin{center}\rule{0.5\linewidth}{0.5pt}\end{center}

\section{13. The Kodaira--Sato Theory}\label{ch1954kunihikokodaira:the-kodairasato-theory}

\subsection{13.1 The Kodaira--Sato
Conjecture}\label{ch1954kunihikokodaira:the-kodairasato-conjecture}

The \textbf{Kodaira--Sato conjecture} (more precisely, the work of
Kodaira and Mikio Sato on the classification of complex analytic
surfaces) concerns the structure of complex surfaces. It states that
every compact complex surface with \(b_1\) odd is of class VII.

\subsection{13.2 Class VII Surfaces}\label{ch1954kunihikokodaira:class-vii-surfaces}

\textbf{Class VII surfaces} are non-Kähler compact complex surfaces with
\(\kappa = -\infty\) and \(b_1 = 1\). They include: - \textbf{Hopf
surfaces:} Quotients of \(\mathbb{C}^2 \setminus \{0\}\) by a
contraction. - \textbf{Inoue surfaces:} Certain surfaces with no curves
at all. - \textbf{Kato surfaces:} A more exotic class.

\begin{center}\rule{0.5\linewidth}{0.5pt}\end{center}

\section{14. The Lefschetz Theorem on
(1,1)-Classes}\label{ch1954kunihikokodaira:the-lefschetz-theorem-on-11-classes}

\subsection{14.1 Statement}\label{ch1954kunihikokodaira:statement-3}

The \textbf{Lefschetz theorem on (1,1)-classes} (proved by Kodaira and
Spencer using Hodge theory) states:

\begin{quote}
\textbf{Lefschetz (1,1) Theorem:} On a compact Kähler manifold \(X\), a
cohomology class in \(H^2(X, \mathbb{C})\) is of type \((1,1)\) if and
only if it lies in the image of the natural map
\(H^2(X, \mathbb{Z}) \to H^2(X, \mathbb{C})\).
\end{quote}

\subsection{14.2 The Hodge Conjecture}\label{ch1954kunihikokodaira:the-hodge-conjecture}

The Lefschetz (1,1) theorem is the only fully proven case of the
\textbf{Hodge conjecture}. For classes of type \((p,p)\) with \(p > 1\),
the Hodge conjecture remains one of the important open problems in
mathematics.

\begin{center}\rule{0.5\linewidth}{0.5pt}\end{center}

\section{15. The Kodaira--Spencer Approach to
Moduli}\label{ch1954kunihikokodaira:the-kodairaspencer-approach-to-moduli}

\subsection{15.1 Moduli of Curves}\label{ch1954kunihikokodaira:moduli-of-curves}

Kodaira and Spencer developed a general approach to moduli problems
using deformation theory. For a curve \(C\) of genus \(g \geq 2\), they
showed: - The moduli space \(\mathcal{M}_g\) has dimension \(3g - 3\). -
The tangent space to \(\mathcal{M}_g\) at \([C]\) is
\(H^1(C, \Theta_C)\). - Deformations are unobstructed because
\(H^2(C, \Theta_C) = 0\).

\subsection{15.2 The Period Map}\label{ch1954kunihikokodaira:the-period-map}

The \textbf{period map} sends a complex structure on a Kähler manifold
to the Hodge decomposition of its cohomology. For a family
\(\pi: \mathcal{X} \to B\), the period map is:

\[P: B \to D / \Gamma\]

where \(D\) is a period domain parametrizing Hodge structures and
\(\Gamma\) is the monodromy group. Kodaira's work on harmonic integrals
provided the analytic foundation for the study of the period map, which
was later developed by Griffiths.

\subsection{15.3 The Torelli Problem}\label{ch1954kunihikokodaira:the-torelli-problem}

The \textbf{Torelli problem} asks whether a variety is determined by its
Hodge structure. Kodaira's embedding theorem and the Hodge decomposition
give partial results, but the full Torelli problem remains unsolved for
most classes of varieties.

\subsection{15.4 Moduli of Higher-Dimensional
Manifolds}\label{ch1954kunihikokodaira:moduli-of-higher-dimensional-manifolds}

For higher-dimensional manifolds, the moduli problem is more complicated
due to potential obstructions. The Kodaira--Spencer theory provides the
local structure of the moduli space in terms of cohomology groups. The
\textbf{global moduli space} is constructed using the theory of stacks
or analytic spaces.

\begin{center}\rule{0.5\linewidth}{0.5pt}\end{center}

\section{16. The Kodaira--Spencer--Kuranishi
Theory}\label{ch1954kunihikokodaira:the-kodairaspencerkuranishi-theory}

\subsection{16.1 The Kuranishi Family}\label{ch1954kunihikokodaira:the-kuranishi-family}

The \textbf{Kuranishi family} is a versal deformation of a compact
complex manifold:

\begin{quote}
\textbf{Kuranishi's Theorem:} There exists a finite-dimensional analytic
space \(B\) and a family \(\pi: \mathcal{X} \to B\) such that every
deformation of \(X_0\) is obtained by pullback, and the tangent space of
\(B\) at \(0\) is \(H^1(X_0, \Theta_{X_0})\).
\end{quote}

\subsection{16.2 Kodaira and Spencer's
Contributions}\label{ch1954kunihikokodaira:kodaira-and-spencers-contributions}

Kodaira and Spencer developed the analytic foundations of deformation
theory, including: - The theory of \textbf{differentiable families} of
complex structures. - The \textbf{upper semicontinuity theorem} for
cohomology groups. - The \textbf{Grauert's direct image theorem} (proved
by Grauert based on Kodaira--Spencer ideas).

\begin{center}\rule{0.5\linewidth}{0.5pt}\end{center}

\section{17. The Theory of Pluricanonical
Maps}\label{ch1954kunihikokodaira:the-theory-of-pluricanonical-maps}

\subsection{17.1 Pluricanonical Systems}\label{ch1954kunihikokodaira:pluricanonical-systems}

The \(m\)-th \textbf{pluricanonical map} is the rational map:

\[\varphi_{mK}: X \dashrightarrow \mathbb{P}(H^0(X, K_X^m)^*)\]

defined by the complete linear system \(|mK_X|\). Kodaira's vanishing
theorem ensures that for sufficiently large \(m\), this map is
well-defined and gives a birational embedding when
\(\kappa(X) = \dim X\).

\subsection{17.2 The Kodaira--Enriques
Classification}\label{ch1954kunihikokodaira:the-kodairaenriques-classification}

For surfaces, the pluricanonical maps are the basis of the
classification: - For \(\kappa = -\infty\), all plurigenera vanish, and
the surface is either rational, ruled, or of class VII. - For
\(\kappa = 0\), \(P_m = 0\) or \(1\), and the \(m\)-canonical map is
finite onto a variety of dimension \(0\). - For \(\kappa = 1\),
\(P_m \sim cm\), and the \(m\)-canonical map fibers the surface over a
curve (elliptic surfaces). - For \(\kappa = 2\), \(P_m \sim cm^2\), and
the \(m\)-canonical map is birational onto a surface of general type.

\subsection{17.3 The Geography of
Surfaces}\label{ch1954kunihikokodaira:the-geography-of-surfaces}

The \textbf{geography of surfaces of general type} studies the possible
values of the invariants \(c_1^2\) (the self-intersection of the
canonical class) and \(c_2\) (the Euler characteristic). The
\textbf{Noether inequality} states:

\[c_1^2 \geq 2p_g - 4\]

where \(p_g = \dim H^0(X, K_X)\) is the geometric genus. The
\textbf{Castelnuovo inequality} and the \textbf{Bogomolov--Miyaoka--Yau
inequality} give further restrictions.

\begin{center}\rule{0.5\linewidth}{0.5pt}\end{center}

\section{18. The Kodaira--Bogomolov
Theorem}\label{ch1954kunihikokodaira:the-kodairabogomolov-theorem}

\subsection{18.1 Statement}\label{ch1954kunihikokodaira:statement-4}

The \textbf{Kodaira--Bogomolov theorem} (proved by Bogomolov building on
Kodaira's ideas) states that a compact complex manifold with
\(c_1(X) = 0\) (i.e., with torsion canonical bundle) has a finite cover
that is a product of a Calabi--Yau manifold, a complex torus, and a
holomorphic symplectic manifold.

\subsection{18.2 The Beauville--Bogomolov
Decomposition}\label{ch1954kunihikokodaira:the-beauvillebogomolov-decomposition}

This leads to the \textbf{Beauville--Bogomolov decomposition theorem}:
every compact Kähler manifold with \(c_1 = 0\) has a finite étale cover
that is a product of: - \textbf{Calabi--Yau manifolds} (holonomy
\(SU(n)\)). - \textbf{Irreducible holomorphic symplectic manifolds}
(holonomy \(Sp(n)\)). - \textbf{Complex tori} (holonomy \(U(n)\) with
flat metric).

\begin{center}\rule{0.5\linewidth}{0.5pt}\end{center}

\section{19. The Kodaira--Néron Classification of Singular
Fibers}\label{ch1954kunihikokodaira:the-kodairanuxe9ron-classification-of-singular-fibers}

\subsection{19.1 The Néron Model}\label{ch1954kunihikokodaira:the-nuxe9ron-model}

The \textbf{Kodaira--Néron classification} (developed by Kodaira and
André Néron) classifies the possible singular fibers of an elliptic
fibration. The classification is in terms of the \textbf{Kodaira symbol}
and the \textbf{Néron type}:

{\def\LTcaptype{none} % do not increment counter
\begin{longtable}[]{@{}
 >{\raggedright\arraybackslash}p{(\linewidth - 4\tabcolsep) * \real{0.3333}}
 >{\raggedright\arraybackslash}p{(\linewidth - 4\tabcolsep) * \real{0.3333}}
 >{\raggedright\arraybackslash}p{(\linewidth - 4\tabcolsep) * \real{0.3333}}@{}}
\toprule\noalign{}
\begin{minipage}[b]{\linewidth}\raggedright
Kodaira symbol
\end{minipage} & \begin{minipage}[b]{\linewidth}\raggedright
Fiber type
\end{minipage} & \begin{minipage}[b]{\linewidth}\raggedright
Description
\end{minipage} \\
\midrule\noalign{}
\endhead
\bottomrule\noalign{}
\endlastfoot
\(I_0\) & Smooth & Smooth elliptic curve \\
\(I_1\) & Nodal & Rational curve with one node \\
\(I_2\) & Two rational curves & Meeting transversely at two points \\
\(I_n\) (\(n \geq 3\)) & Cycle of \(n\) rational curves & Cyclic
configuration \\
\(II\) & Cuspidal & Rational curve with a cusp \\
\(III\) & Two rational curves & Meeting tangentially \\
\(IV\) & Three rational curves & Meeting at a common point \\
\end{longtable}
}

\subsection{19.2 The Tate Algorithm}\label{ch1954kunihikokodaira:the-tate-algorithm}

The \textbf{Tate algorithm} provides a practical method for determining
the Kodaira--Néron type of a singular fiber of an elliptic curve over a
local field. By analyzing the Weierstrass equation of the elliptic
curve, one can compute the Kodaira type from the valuations of the
coefficients.

\subsection{19.3 Applications}\label{ch1954kunihikokodaira:applications-2}

This classification is essential for: - The study of \textbf{elliptic
surfaces} and their moduli. - \textbf{Arithmetic geometry}: reduction
types of elliptic curves. - \textbf{F-theory}: the physics of string
theory compactifications. - \textbf{Desingularization} of surfaces with
elliptic fibrations.

\begin{center}\rule{0.5\linewidth}{0.5pt}\end{center}

\section{20. The Kodaira--Spencer Stability
Theorem}\label{ch1954kunihikokodaira:the-kodairaspencer-stability-theorem}

\subsection{20.1 Stability of Complex
Structures}\label{ch1954kunihikokodaira:stability-of-complex-structures}

The \textbf{Kodaira--Spencer stability theorem} states that a compact
complex manifold with \(H^1(X, \Theta_X) = 0\) is \textbf{rigid}: it
admits no non-trivial deformations.

\subsection{20.2 Examples}\label{ch1954kunihikokodaira:examples}

\begin{itemize}
\tightlist
\item
 The projective plane \(\mathbb{P}^n\) is rigid:
 \(H^1(\mathbb{P}^n, \Theta_{\mathbb{P}^n}) = 0\).
\item
 A generic hypersurface of degree \(d \geq 3\) in \(\mathbb{P}^n\) is
 rigid when \(n \geq 5\).
\item
 Calabi--Yau manifolds typically have non-trivial moduli spaces, with
 dimension \(h^{1, n-1}\).
\end{itemize}

\subsection{20.3 The Kodaira--Spencer Map and
Obstructions}\label{ch1954kunihikokodaira:the-kodairaspencer-map-and-obstructions}

The \textbf{obstruction theory} for deformations is encoded in the
\textbf{Kodaira--Spencer obstruction classes}. The first obstruction
lies in \(H^2(X, \Theta_X)\). If this group vanishes, deformations are
unobstructed and the moduli space is smooth of dimension
\(h^1(X, \Theta_X)\).

\begin{center}\rule{0.5\linewidth}{0.5pt}\end{center}

\section{21. Kodaira's Influence on Modern Algebraic
Geometry}\label{ch1954kunihikokodaira:kodairas-influence-on-modern-algebraic-geometry}

\subsection{21.1 The Minimal Model
Program}\label{ch1954kunihikokodaira:the-minimal-model-program}

Kodaira's work on the classification of surfaces inspired the
\textbf{minimal model program} (MMP) in higher dimensions. The MMP aims
to classify all projective varieties by: 1. Constructing a minimal model
with \(K_X\) nef (or a Mori fiber space with \(K_X\) negative). 2. Using
the Kodaira dimension as the primary invariant. 3. Applying vanishing
theorems (generalizations of Kodaira's) to control cohomology.

\subsection{21.2 The Hodge Conjecture}\label{ch1954kunihikokodaira:the-hodge-conjecture-1}

Kodaira's work on the Lefschetz (1,1) theorem and Hodge theory provides
the only fully solved case of the \textbf{Hodge conjecture}. The general
Hodge conjecture --- that every Hodge class is algebraic --- remains
open.

\subsection{21.3 Mirror Symmetry}\label{ch1954kunihikokodaira:mirror-symmetry}

Kodaira's deformation theory and classification of surfaces are
essential tools in \textbf{mirror symmetry}, the conjectural duality
between Calabi--Yau manifolds. The moduli spaces of complex structures
on a Calabi--Yau manifold are mirrored by the moduli spaces of Kähler
structures on its mirror.

\begin{center}\rule{0.5\linewidth}{0.5pt}\end{center}

\section{22. Awards and Recognition}\label{ch1954kunihikokodaira:awards-and-recognition}

\textbf{1954:} - Fields Medal (International Mathematical Union) ---
awarded at the ICM in Amsterdam.

\textbf{1957:} - Japan Academy Prize --- for outstanding contributions
to mathematics. - Order of Culture (Bunka Kōrōsha) --- the highest
Japanese honor in culture.

\textbf{1975:} - Fujiwara Prize.

\textbf{1984/85:} - Wolf Prize in Mathematics --- ``For his outstanding
contributions to the study of complex manifolds and algebraic
varieties.''

\textbf{Elected to:} - Japan Academy. - American Academy of Arts and
Sciences. - National Academy of Sciences (USA).

\begin{center}\rule{0.5\linewidth}{0.5pt}\end{center}

\section{22. Selected Major
Publications}\label{ch1954kunihikokodaira:selected-major-publications}

\begin{enumerate}
\def\labelenumi{\arabic{enumi}.}
\item
 \textbf{``The theory of harmonic integrals and their applications to
 algebraic geometry.''} (with D. C. Spencer) \emph{Princeton University
 Press}, 1953. --- \emph{The foundational text on the applications of
 Hodge theory.}
\item
 \textbf{``On Kähler varieties of restricted type.''} \emph{Annals of
 Mathematics} 60, 28--48, 1954. --- \emph{The Kodaira embedding
 theorem.}
\item
 \textbf{``On the structure of compact complex analytic surfaces.''} (4
 papers) \emph{American Journal of Mathematics} 86--88, 1964--1966. ---
 \emph{The classification of complex surfaces.}
\item
 \textbf{``On deformations of complex analytic structures.''} (with D.
 C. Spencer) \emph{Annals of Mathematics} 67, 328--466, 1958. ---
 \emph{The foundational paper on deformation theory.}
\item
 \textbf{``A theorem of completeness of characteristic systems for
 analytic families of compact submanifolds of complex manifolds.''}
 \emph{Annals of Mathematics} 75, 146--162, 1962. --- \emph{The
 completeness of the characteristic system.}
\item
 \textbf{``On the structure of compact complex analytic surfaces,
 I--IV.''} \emph{American Journal of Mathematics} 86(4), 751--798,
 1964; 88(3), 682--721, 1966. --- \emph{The complete classification.}
\item
 \textbf{``Harmonic fields in Riemannian manifolds.''} \emph{Annals of
 Mathematics} 50, 587--665, 1949. --- \emph{Generalized Hodge theory.}
\item
 \textbf{``Complex manifolds and deformation of complex structures.''}
 Springer, 1986. --- \emph{The definitive monograph on deformation
 theory.}
\item
 \textbf{``Pluricanonical systems on algebraic surfaces.''}
 \emph{Nagoya Mathematical Journal} 34, 1--21, 1969. ---
 \emph{Pluricanonical maps and the classification.}
\end{enumerate}

\begin{center}\rule{0.5\linewidth}{0.5pt}\end{center}

\section{23. References and Further
Reading}\label{ch1954kunihikokodaira:references-and-further-reading}

\subsection{Official Sources}\label{ch1954kunihikokodaira:official-sources}

\begin{itemize}
\tightlist
\item
 IMU Fields Medal 1954 Announcement:
 https://www.mathunion.org/imu-awards/fields-medal/fields-medals-1954
\item
 Wolf Prize Citation, 1984.
\item
 Japan Academy: Kunihiko Kodaira Biography.
\end{itemize}

\subsection{Key Technical References}\label{ch1954kunihikokodaira:key-technical-references}

\begin{itemize}
\tightlist
\item
 Kodaira, K. \emph{Collected Works} (3 vols.). Iwanami Shoten and
 Princeton University Press, 1975.
\item
 Kodaira, K. \emph{Complex Manifolds and Deformation of Complex
 Structures}. Springer, 1986.
\item
 Morrow, J. and Kodaira, K. \emph{Complex Manifolds}. Holt, Rinehart
 and Winston, 1971.
\end{itemize}

\subsection{Textbooks and Surveys}\label{ch1954kunihikokodaira:textbooks-and-surveys}

\begin{itemize}
\tightlist
\item
 Griffiths, P. and Harris, J. \emph{Principles of Algebraic Geometry}.
 Wiley, 1978.
\item
 Barth, W., Peters, C., and Van de Ven, A. \emph{Compact Complex
 Surfaces}. Springer, 1984.
\item
 Beauville, A. \emph{Complex Algebraic Surfaces}. Cambridge University
 Press, 1996.
\item
 Huybrechts, D. \emph{Complex Geometry}. Springer, 2005.
\item
 Friedman, R. \emph{Algebraic Surfaces and Holomorphic Vector Bundles}.
 Springer, 1998.
\item
 Kollár, J. and Mori, S. \emph{Birational Geometry of Algebraic
 Varieties}. Cambridge University Press, 1998.
\end{itemize}

\subsection{Biography and Exposition}\label{ch1954kunihikokodaira:biography-and-exposition}

\begin{itemize}
\tightlist
\item
 Weyl, H. ``The Work of Kunihiko Kodaira.'' ICM 1954 Laudatio.
\item
 Spencer, D. C. ``Kunihiko Kodaira.'' \emph{Notices of the AMS} 45,
 1998.
\item
 Iyanaga, S. ``Kunihiko Kodaira: A Reminiscence.'' \emph{Mathematical
 Intelligencer} 21, 1999.
\end{itemize}

\begin{center}\rule{0.5\linewidth}{0.5pt}\end{center}

\emph{(End of report)}
